\documentclass[12pt,a4paper]{article}
\usepackage[utf8]{inputenc}
\usepackage[T5]{fontenc}
\usepackage[vietnamese]{babel}
\usepackage[margin=1.5cm,top=1.2cm,bottom=1.2cm]{geometry}
\usepackage{amsmath,amssymb,array,tabularx,enumitem,fancyhdr,tikz,xcolor}
\usetikzlibrary{calc,angles,quotes,arrows.meta,babel,patterns}
\definecolor{ink}{RGB}{30,55,70}

\newcommand{\headerQuote}{%
  \ifcase\value{page}\relax
    \or Math gives us hope that every problem has a solution% Trang 1
    \or Math + You = Champion% Trang 2
    \or Where Frustration Becomes Inspiration% Trang 3
    \or Turning Confusion into Clarity% Trang 4
    \or Don't be an X, be a Y of mathematics% Trang 5
    \or Don't get mad. Learn to add% Trang 6
    \or Math gives us hope that every problem has a solution% Trang 7
    \or Math + You = Champion% Trang 8
    \or Where Frustration Becomes Inspiration% Trang 9
    \or Turning Confusion into Clarity% Trang 10
    \or Don't be an X, be a Y of mathematics% Trang 11
    \or Don't get mad. Learn to add% Trang 12
    \else Math gives us hope that every problem has a solution%
  \fi
}
\pagestyle{fancy}\fancyhf{}
\lhead{\small\textit{\headerQuote}}
\rhead{\small\textit{Người biên soạn: Dương Gia Bảo}}
\rfoot{\small Trang \thepage}
\renewcommand{\headrulewidth}{0.4pt}
\setlength{\headheight}{15pt}
\setlength{\parindent}{0pt}
\setlength{\parskip}{0.8pt}
\setlength{\tabcolsep}{2.5pt}
\renewcommand{\arraystretch}{1.08}
\newcommand{\parttitle}[1]{\par\vspace{1.5pt}{\large\bfseries\color{ink}#1}\par\vspace{1pt}}
\newcommand{\four}[4]{\par\noindent\begin{tabularx}{\linewidth}{@{}*{4}{>{\raggedright\arraybackslash}X}@{}}\textbf{A.} #1&\textbf{B.} #2&\textbf{C.} #3&\textbf{D.} #4\end{tabularx}}
\newcommand{\two}[4]{\par\noindent\begin{tabularx}{\linewidth}{@{}*{2}{>{\raggedright\arraybackslash}X}@{}}\textbf{A.} #1&\textbf{B.} #2\\[1.2pt]\textbf{C.} #3&\textbf{D.} #4\end{tabularx}}
\newcommand{\answerbox}{\hfill\textit{Đáp số:}\ \fbox{\rule{0pt}{10pt}\hspace{2.2cm}}}
\setlist[enumerate,1]{label=\textbf{Câu \arabic*.},leftmargin=1.35cm,labelsep=.15cm,itemsep=0.8pt,topsep=0.8pt,parsep=0pt}
\setlist[enumerate,2]{label=\alph*),leftmargin=.55cm,itemsep=0.5pt,topsep=0.5pt,parsep=0pt}

\begin{document}
\noindent
\begin{minipage}[c]{0.38\textwidth}
    \centering
    \textbf{LÀM TOÁN ĐỂ HẠNH PHÚC}\\[-2pt]
    \rule{3.2cm}{0.6pt}\\[2pt]
    \textbf{ĐỀ CHÍNH THỨC}
\end{minipage}\hfill
\begin{minipage}[c]{0.35\textwidth}
    \centering
    \textbf{ĐỀ ÔN THI GK1}\\
    \textbf{NĂM HỌC 2026 -- 2027}\\[1pt]
    \textbf{ĐỀ SỐ 16}\\
    \textbf{Môn: TOÁN -- Khối lớp 10}\\
    \textit{\small Thời gian làm bài: 90 phút\\(không kể thời gian phát đề)}
\end{minipage}

\vspace{0.1cm}
\noindent
\begin{minipage}[b]{0.78\textwidth}
    \textbf{Họ và tên:}\dotfill\textbf{SBD:}\makebox[3.5cm]{\dotfill}
\end{minipage}\hfill
\begin{minipage}[b]{0.20\textwidth}
    \hfill
    \setlength{\fboxrule}{0.8pt}%
    \fbox{\begin{minipage}{2.3cm}
        \centering
        \textbf{\small Mã đề thi}\\
        \textbf{0016}
    \end{minipage}}
\end{minipage}
\vspace{0.1cm}

\parttitle{PHẦN I. Trắc nghiệm nhiều lựa chọn}
\textit{Thí sinh trả lời từ câu 1 đến câu 12. Mỗi câu hỏi chỉ chọn một phương án đúng duy nhất.}
\begin{enumerate}
\item Cho tập hợp $A = \{a; b; c\}$. Khẳng định nào sau đây là \textbf{đúng}?
\four{$\{a\} \subset A$.}{$a \subset A$.}{$\{a; b\} \in A$.}{$\varnothing \in A$.}

\item Cho định lý: ``Nếu hai tam giác bằng nhau thì chúng có diện tích bằng nhau''. Khẳng định nào sau đây là \textbf{đúng}?
\two{Diện tích bằng nhau là điều kiện đủ để hai tam giác bằng nhau.}{Hai tam giác bằng nhau là điều kiện đủ để diện tích bằng nhau.}{Hai tam giác bằng nhau là điều kiện cần để diện tích bằng nhau.}{Diện tích bằng nhau là điều kiện cần và đủ để hai tam giác bằng nhau.}

\item Trong các mệnh đề sau đây, mệnh đề nào là một mệnh đề \textbf{đúng}?
\two{Số $15$ là một số nguyên tố.}{Tổng hai góc nhọn của một tam giác vuông bằng $180^\circ$.}{Tứ giác có bốn cạnh bằng nhau là hình thoi.}{Phương trình $x^2 + 1 = 0$ có nghiệm thực.}

\item Cho tập hợp $X = \{1; 2; 3; 4\}$. Số tập hợp con gồm đúng $2$ phần tử của tập hợp $X$ là
\four{$4$.}{$8$.}{$16$.}{$6$.}

\item Cặp số $(x; y)$ nào sau đây là một nghiệm của hệ bất phương trình $\begin{cases} 2x - y + 1 > 0 \\ x + 3y - 6 \le 0 \end{cases}$?
\four{$(1; 1)$.}{$(0; 3)$.}{$(-2; 1)$.}{$(4; 2)$.}

\item Hệ bất phương trình nào sau đây là một hệ bất phương trình bậc nhất hai ẩn?
\four{$\begin{cases} 2x^2 + y \le 4 \\ x - 3y > 1 \end{cases}$}{$\begin{cases} 3x - y + 2 \ge 0 \\ 2x + 5y - 7 < 0 \end{cases}$}{$\begin{cases} x - y + z \le 0 \\ 2x + y > 3 \end{cases}$}{$\begin{cases} xy + 2 \le 0 \\ x - 4y > 0 \end{cases}$}

\item Miền nghiệm của bất phương trình $3x - 2y \ge 6$ chứa điểm nào sau đây?
\four{$O(0; 0)$.}{$M(1; 2)$.}{$N(4; 1)$.}{$P(2; 3)$.}

\item Bất phương trình nào sau đây là một bất phương trình bậc nhất hai ẩn?
\four{$2x^2 - 3y + 1 > 0$.}{$x + 2y - 3z \le 0$.}{$3x + \dfrac{2}{y} - 5 < 0$.}{$4x - 5y + 1 \ge 0$.}

\item Cho tam giác $ABC$ có các cạnh $BC = a, CA = b, AB = c$. Gọi $m_a$ là độ dài đường trung tuyến kẻ từ đỉnh $A$. Khẳng định nào sau đây là \textbf{đúng}?
\four{$m_a^2 = \tfrac{2(b^2 + c^2) - a^2}{4}$.}{$m_a^2 = \tfrac{b^2 + c^2 - 2a^2}{4}$.}{$m_a^2 = \tfrac{2(b^2 + c^2) + a^2}{4}$.}{$m_a^2 = \tfrac{b^2 + c^2 - a^2}{2}$.}

\item Giá trị của $\sin 150^\circ$ bằng
\four{$-\tfrac{1}{2}$.}{$\tfrac{1}{2}$.}{$\tfrac{\sqrt{3}}{2}$.}{$-\tfrac{\sqrt{3}}{2}$.}

\item Với mọi góc $\alpha$ thỏa mãn $0^\circ \le \alpha \le 180^\circ$, khẳng định nào sau đây là \textbf{đúng}?
\four{$\sin(180^\circ - \alpha) = -\sin \alpha$.}{$\cos(180^\circ - \alpha) = \cos \alpha$.}{$\cos(180^\circ - \alpha) = -\cos \alpha$.}{$\tan(180^\circ - \alpha) = \tan \alpha$.}

\item Cho góc tù $\alpha$ ($90^\circ < \alpha < 180^\circ$). Khẳng định nào sau đây là \textbf{đúng}?
\four{$\sin \alpha < 0$.}{$\cos \alpha > 0$.}{$\tan \alpha > 0$.}{$\cot \alpha < 0$.}
\end{enumerate}

\parttitle{PHẦN II. Trắc nghiệm đúng/sai}
\textit{Thí sinh trả lời từ câu 13 đến câu 14. Trong mỗi ý a), b), c), d) ở mỗi câu, thí sinh chọn đúng hoặc sai.}
\begin{enumerate}[resume]
\item Cho hai tập hợp con của tập số thực $\mathbb{R}$ là $A = [-4; 3)$ và $B = (0; 5]$.
\begin{enumerate}
    \item Số $0$ thuộc tập hợp $A$ nhưng không thuộc tập hợp $B$.
    \item Tập hợp $B$ là nửa khoảng gồm tất cả các số thực $x$ thỏa mãn $0 < x \le 5$.
    \item Giao của hai tập hợp là khoảng $A \cap B = (0; 3)$.
    \item Hiệu của hai tập hợp là nửa khoảng $A \setminus B = [-4; 0)$.
\end{enumerate}

\item Xét góc nhọn $\alpha$ bất kì ($0^\circ < \alpha < 90^\circ$) và đặt $\beta = 180^\circ - \alpha$.
\begin{enumerate}
    \item Hai góc $\alpha$ và $\beta$ là hai góc bù nhau.
    \item Ta có đẳng thức $\sin \beta = \sin \alpha$.
    \item Biểu thức $M = \sin^2 \alpha + \sin^2(90^\circ - \alpha)$ có giá trị không đổi và bằng $1$.
    \item Khi $\alpha = 45^\circ$, giá trị của biểu thức $N = \cos \alpha + \cos \beta + \tan \alpha + \cot \beta$ bằng $2$.
\end{enumerate}
\end{enumerate}

\parttitle{PHẦN III. Trắc nghiệm trả lời ngắn}
\textit{Thí sinh trả lời từ câu 15 đến câu 18. Viết kết quả vào ô quy định.}
\begin{enumerate}[resume]
\item Cho hai tập hợp số $A = [-3; 7)$ và $B = (1; 10]$. Tìm số phần tử nguyên thuộc tập hợp giao $C = A \cap B$.
\answerbox

\item Trong mặt phẳng tọa độ $Oxy$, cho tam giác $ABC$ có ba đỉnh là $A(1; 2)$, $B(4; 1)$ và $C(2; 5)$. Miền tam giác $ABC$ (kể cả biên) là miền nghiệm của một hệ bất phương trình bậc nhất hai ẩn. Tìm giá trị nguyên nhỏ nhất của tham số $m$ để bất phương trình $2x + 3y \le m$ nghiệm đúng với mọi điểm $(x; y)$ thuộc miền tam giác $ABC$.
\answerbox

\item Cho tam giác $ABC$ có độ dài ba cạnh thỏa mãn hệ thức: $b^2 + c^2 - a^2 = \sqrt{3}bc$. Tính số đo của góc $A$ (theo đơn vị độ).
\answerbox

\item Cho hai tập hợp hữu hạn $A = \{x \in \mathbb{N} \mid (x - 2)(x^2 - 5x + 6) = 0\}$ và $B = \{x \in \mathbb{Z} \mid -1 \le x \le 4\}$. Tìm số phần tử của tập hợp $B \setminus A$.
\answerbox
\end{enumerate}

\parttitle{PHẦN IV. Tự luận}
\textit{Thí sinh trình bày lời giải chi tiết từ câu 19 đến câu 22.}
\begin{enumerate}[resume]
\item \textit{(1,0 điểm)} Một khu đất công viên có dạng hình tam giác $ABC$ như hình vẽ bên. Ban quản lý đo đạc được độ dài hai cạnh là $AB = 50\text{ m}$, $AC = 80\text{ m}$ và góc xen giữa $\widehat{BAC} = 60^\circ$.
\begin{minipage}[t]{0.62\linewidth}
\begin{enumerate}
    \item Tính độ dài cạnh $BC$ của khu đất công viên.
    \item Tính diện tích $S$ của khu đất công viên (lấy giá trị đúng theo $\sqrt{3}$).
    \item Ban quản lý dự định xây dựng một hồ nước hình tròn nội tiếp trong tam giác $ABC$. Tính bán kính $r$ của hồ nước này (\textit{kết quả làm tròn đến hàng phần mười}).
\end{enumerate}
\end{minipage}%
\hfill
\begin{minipage}[t]{0.36\linewidth}
\centering
\begin{tikzpicture}[baseline=(current bounding box.north), scale=0.82, >=Stealth]
    % Điểm A(0,0), C(3.2, 0), B(1.0, 1.732)
    \coordinate (A) at (0,0);
    \coordinate (C) at (3.2,0);
    \coordinate (B) at (1.0, 1.732);
    \coordinate (I) at (1.2, 0.693);

    % Miền tam giác
    \fill[fill=teal!12] (A) -- (B) -- (C) -- cycle;
    \draw[line width=1.2pt, ink] (A) -- (B) -- (C) -- cycle;

    % Hồ nước hình tròn nội tiếp
    \fill[fill=cyan!35] (I) circle (0.693);
    \draw[line width=1.1pt, blue!80!black] (I) circle (0.693);
    \filldraw[black] (I) circle (1.8pt) node[right, font=\footnotesize\bfseries] {$I$};

    % Điểm đỉnh
    \filldraw[black] (A) circle (2pt) node[below left, font=\small\bfseries] {$A$};
    \filldraw[black] (B) circle (2pt) node[above, font=\small\bfseries] {$B$};
    \filldraw[black] (C) circle (2pt) node[below right, font=\small\bfseries] {$C$};

    % Góc 60 độ
    \draw[line width=0.8pt] (0.45,0) arc (0:60:0.45);
    \node at (0.65, 0.28) {\footnotesize $60^\circ$};

    % Nhãn kích thước
    \node[left, font=\footnotesize] at (0.45, 1.0) {$50\text{ m}$};
    \node[below, font=\footnotesize] at (1.6, -0.05) {$80\text{ m}$};
    \node[above right, font=\footnotesize, teal!80!black] at (2.1, 0.95) {$BC = ?$};
\end{tikzpicture}
\end{minipage}

\item \textit{(0,75 điểm)} Để chuẩn bị cho đại hội thể thao cấp trường, Đoàn trường tổ chức cuộc khảo sát đối với $45$ học sinh của lớp 10A về việc đăng ký tham gia ba môn thi đấu: Bóng đá, Cầu lông và Bóng rổ. Kết quả thu được như sau:
\begin{itemize}[leftmargin=0.5cm, itemsep=0.5pt]
    \item Có $25$ học sinh đăng ký thi đấu Bóng đá; $20$ học sinh đăng ký Cầu lông; $15$ học sinh đăng ký Bóng rổ;
    \item Có $10$ học sinh đăng ký cả Bóng đá và Cầu lông; $7$ học sinh đăng ký cả Cầu lông và Bóng rổ; $8$ học sinh đăng ký cả Bóng đá và Bóng rổ;
    \item Có $4$ học sinh đăng ký tham gia cả ba môn thể thao trên.
\end{itemize}
Hỏi trong lớp 10A:
\begin{enumerate}
    \item Có bao nhiêu học sinh đăng ký tham gia ít nhất một trong ba môn thể thao nói trên?
    \item Có bao nhiêu học sinh không đăng ký tham gia bất kì môn thể thao nào?
    \item Có bao nhiêu học sinh chỉ đăng ký tham gia đúng một môn thể thao?
\end{enumerate}

\item \textit{(0,75 điểm)} Một cơ sở thủ công mỹ nghệ sản xuất hai loại sản phẩm lưu niệm: Đèn lồng truyền thống (loại A) và Hộp gỗ chạm khảm (loại B).
\begin{itemize}[leftmargin=0.5cm, itemsep=0.5pt]
    \item Để sản xuất một chiếc đèn lồng loại A cần $2$ giờ gia công và $1$ giờ hoàn thiện.
    \item Để sản xuất một chiếc hộp gỗ loại B cần $1$ giờ gia công và $2$ giờ hoàn thiện.
\end{itemize}
Mỗi tuần cơ sở có tối đa $100$ giờ gia công và $80$ giờ hoàn thiện. Biết rằng mỗi chiếc đèn lồng loại A mang lại lợi nhuận $120\,000$ đồng, mỗi chiếc hộp gỗ loại B mang lại lợi nhuận $100\,000$ đồng. Cơ sở cam kết tiêu thụ hết toàn bộ sản phẩm làm ra trong tuần. Hỏi cơ sở cần sản xuất bao nhiêu chiếc đèn lồng loại A và bao nhiêu chiếc hộp gỗ loại B mỗi tuần để tổng số lợi nhuận thu được là lớn nhất?

\item \textit{(0,5 điểm)} Một câu lạc bộ ngoại khóa của trường THPT có $50$ học sinh tham gia. Trong đó có $32$ bạn tham gia Câu lạc bộ Tin học, $28$ bạn tham gia Câu lạc bộ Tiếng Anh và $18$ bạn tham gia cả hai câu lạc bộ này.
\begin{enumerate}
    \item Có bao nhiêu bạn học sinh tham gia ít nhất một trong hai câu lạc bộ Tin học hoặc Tiếng Anh?
    \item Có bao nhiêu bạn học sinh không tham gia cả hai câu lạc bộ trên?
    \item Có bao nhiêu bạn học sinh chỉ tham gia Câu lạc bộ Tin học mà không tham gia Câu lạc bộ Tiếng Anh?
\end{enumerate}
\end{enumerate}

\newpage
% =========================================================================
% BẢNG ĐÁP ÁN VÀ LỜI GIẢI CHI TIẾT
% =========================================================================
\parttitle{BẢNG ĐÁP ÁN TỔNG HỢP}

\vspace{0.15cm}
\textbf{1. Phần I: Trắc nghiệm nhiều lựa chọn}
\begin{center}
{\setlength{\tabcolsep}{2.5pt}\small
\begin{tabular}{|c|c|c|c|c|c|c|c|c|c|c|c|}
\hline
\textbf{Câu 1} & \textbf{Câu 2} & \textbf{Câu 3} & \textbf{Câu 4} & \textbf{Câu 5} & \textbf{Câu 6} & \textbf{Câu 7} & \textbf{Câu 8} & \textbf{Câu 9} & \textbf{Câu 10} & \textbf{Câu 11} & \textbf{Câu 12} \\
\hline
\textbf{A} & \textbf{B} & \textbf{C} & \textbf{D} & \textbf{A} & \textbf{B} & \textbf{C} & \textbf{D} & \textbf{A} & \textbf{B} & \textbf{C} & \textbf{D} \\
\hline
\end{tabular}}
\end{center}

\vspace{0.15cm}
\textbf{2. Phần II: Trắc nghiệm đúng/sai}
\begin{center}
\begin{tabular}{|c|c|c|c|c|}
\hline
\textbf{Câu hỏi} & \textbf{Ý a)} & \textbf{Ý b)} & \textbf{Ý c)} & \textbf{Ý d)} \\
\hline
\textbf{Câu 13} & \textbf{Đúng} & \textbf{Đúng} & \textbf{Đúng} & \textbf{Sai} \\
\hline
\textbf{Câu 14} & \textbf{Đúng} & \textbf{Đúng} & \textbf{Đúng} & \textbf{Sai} \\
\hline
\end{tabular}
\end{center}

\vspace{0.15cm}
\textbf{3. Phần III: Trắc nghiệm trả lời ngắn}
\begin{center}
\begin{tabular}{|c|c|c|c|}
\hline
\textbf{Câu 15} & \textbf{Câu 16} & \textbf{Câu 17} & \textbf{Câu 18} \\
\hline
\textbf{5} & \textbf{19} & \textbf{30} & \textbf{4} \\
\hline
\end{tabular}
\end{center}

\vspace{0.15cm}
\textbf{4. Phần IV: Tự luận}
\begin{center}
\begin{tabularx}{\linewidth}{|c|>{\raggedright\arraybackslash}X|}
\hline
\textbf{Câu hỏi} & \textbf{Kết quả tóm tắt} \\
\hline
\textbf{Câu 19} & a) $BC = 70\text{ m}$;\quad b) $S = 1000\sqrt{3}\text{ m}^2$;\quad c) $r = 10\sqrt{3}\text{ m} \approx 17{,}3\text{ m}$. \\
\hline
\textbf{Câu 20} & a) $39$ học sinh;\quad b) $6$ học sinh;\quad c) $22$ học sinh. \\
\hline
\textbf{Câu 21} & Sản xuất \textbf{40 chiếc đèn lồng loại A} và \textbf{20 chiếc hộp gỗ loại B}; lợi nhuận lớn nhất là $6\,800\,000$ đồng. \\
\hline
\textbf{Câu 22} & a) $42$ bạn;\quad b) $8$ bạn;\quad c) $14$ bạn. \\
\hline
\end{tabularx}
\end{center}

\vspace{0.25cm}
% =========================================================================
% HÌNH VẼ BIỂU DIỄN MIỀN NGHIỆM (HỆ BẤT PHƯƠNG TRÌNH)
% =========================================================================
\parttitle{HÌNH VẼ BIỂU DIỄN MIỀN NGHIỆM (HỆ BẤT PHƯƠNG TRÌNH)}

\begin{minipage}[t]{0.48\textwidth}
\centering
\textbf{Hình 1: Miền nghiệm tam giác $ABC$ ở Câu 16}\\[5pt]
\begin{tikzpicture}[scale=0.88, >=Stealth]
    % Trục toạ độ
    \draw[line width=1.2pt, ->] (-0.8,0) -- (5.5,0) node[right, font=\normalsize\bfseries] {$x$};
    \draw[line width=1.2pt, ->] (0,-0.8) -- (0,5.8) node[above, font=\normalsize\bfseries] {$y$};
    \node[below left, font=\normalsize\bfseries] at (0,0) {$O$};

    % Miền tam giác ABC: A(1;2), B(4;1), C(2;5)
    \fill[fill=ink!18] (1,2) -- (4,1) -- (2,5) -- cycle;
    \draw[line width=2.0pt, blue!70!black] (1,2) -- (4,1) -- (2,5) -- cycle;

    % Đỉnh
    \filldraw[black] (1,2) circle (3.5pt) node[below left, font=\normalsize\bfseries] {$A(1;2)$};
    \filldraw[black] (4,1) circle (3.5pt) node[right, font=\normalsize\bfseries] {$B(4;1)$};
    \filldraw[black] (2,5) circle (3.5pt) node[above right, font=\normalsize\bfseries] {$C(2;5)$};

    % Đường dóng
    \draw[dotted, thick] (1,0) node[below, font=\small] {$1$} -- (1,2) -- (0,2) node[left, font=\small] {$2$};
    \draw[dotted, thick] (4,0) node[below, font=\small] {$4$} -- (4,1) -- (0,1) node[left, font=\small] {$1$};
    \draw[dotted, thick] (2,0) node[below, font=\small] {$2$} -- (2,5) -- (0,5) node[left, font=\small] {$5$};
\end{tikzpicture}
\end{minipage}\hfill
\begin{minipage}[t]{0.48\textwidth}
\centering
\textbf{Hình 2: Miền nghiệm tứ giác $OABC$ ở Câu 21}\\[5pt]
\begin{tikzpicture}[x=0.055cm, y=0.055cm, >=Stealth]
    % Trục toạ độ
    \draw[line width=1.2pt, ->] (-8,0) -- (68,0) node[right, font=\normalsize\bfseries] {$x$};
    \draw[line width=1.2pt, ->] (0,-8) -- (0,60) node[above, font=\normalsize\bfseries] {$y$};
    \node[below left, font=\normalsize\bfseries] at (0,0) {$O$};

    % Miền tứ giác O(0;0), A(50;0), B(40;20), C(0;40)
    \fill[fill=ink!18] (0,0) -- (50,0) -- (40,20) -- (0,40) -- cycle;
    \draw[line width=2.0pt, blue!70!black] (0,0) -- (50,0) -- (40,20) -- (0,40) -- cycle;

    % Đường biên d1: 2x + y = 100
    \draw[line width=1.3pt, dashed, red!80!black] (32, 36) -- (58, -16) node[below right, font=\small] {$d_1: 2x+y=100$};
    % Đường biên d2: x + 2y = 80
    \draw[line width=1.3pt, dashed, teal!80!black] (-6, 43) -- (60, 10) node[above right, font=\small] {$d_2: x+2y=80$};

    % Đỉnh
    \filldraw[black] (0,0) circle (3.5pt);
    \filldraw[black] (50,0) circle (4.0pt) node[below, font=\normalsize\bfseries] {$A(50;0)$};
    \filldraw[black] (40,20) circle (4.0pt) node[above right, font=\normalsize\bfseries] {$B(40;20)$};
    \filldraw[black] (0,40) circle (4.0pt) node[left, font=\normalsize\bfseries] {$C(0;40)$};

    % Đường dóng tọa độ B
    \draw[dotted, thick] (40,0) node[below, font=\small] {$40$} -- (40,20) -- (0,20) node[left, font=\small] {$20$};
\end{tikzpicture}
\end{minipage}

\newpage
% =========================================================================
% HƯỚNG DẪN GIẢI CHI TIẾT
% =========================================================================
\parttitle{HƯỚNG DẪN GIẢI CHI TIẾT}

\textbf{PHẦN I. Trắc nghiệm nhiều lựa chọn}
\begin{itemize}[leftmargin=0.6cm, itemsep=3pt]
    \item \textbf{Câu 1:} Vì $a \in A$ nên tập hợp $\{a\}$ là một tập hợp con của $A$, viết là $\{a\} \subset A$. Khẳng định A đúng. (Các phương án khác: $a \subset A$ là sai vì $a$ là phần tử; $\{a; b\} \in A$ là sai vì $\{a; b\}$ là tập con; $\varnothing \in A$ là sai vì $\varnothing \subset A$). Chọn \textbf{A}.
    \item \textbf{Câu 2:} Mệnh đề có dạng $P \implies Q$, trong đó $P$ là ``Hai tam giác bằng nhau'' và $Q$ là ``Diện tích bằng nhau''. Theo định nghĩa logic, $P$ là điều kiện đủ để có $Q$. Vậy ``Hai tam giác bằng nhau là điều kiện đủ để diện tích bằng nhau''. Chọn \textbf{B}.
    \item \textbf{Câu 3:} Tứ giác có bốn cạnh bằng nhau là hình thoi (theo định nghĩa hình thoi trong hình học). Đây là mệnh đề đúng. (Số $15$ chia hết cho $3$ và $5$ nên là hợp số; tổng hai góc nhọn trong tam giác vuông bằng $90^\circ$; phương trình $x^2 + 1 = 0$ vô nghiệm thực vì $x^2 + 1 \ge 1 > 0$). Chọn \textbf{C}.
    \item \textbf{Câu 4:} Số tập hợp con gồm $2$ phần tử của tập hợp có $4$ phần tử là số tổ hợp chập $2$ của $4$: $C_4^2 = \dfrac{4!}{2!(4-2)!} = 6$. Các tập hợp con đó là $\{1; 2\}, \{1; 3\}, \{1; 4\}, \{2; 3\}, \{2; 4\}, \{3; 4\}$. Chọn \textbf{D}.
    \item \textbf{Câu 5:} Thay tọa độ điểm $(1; 1)$ vào hệ bất phương trình:
    $\begin{cases} 2(1) - 1 + 1 = 2 > 0 \\ 1 + 3(1) - 6 = -2 \le 0 \end{cases}$ (thỏa mãn cả hai bất phương trình). Do đó $(1; 1)$ là nghiệm của hệ. Chọn \textbf{A}.
    \item \textbf{Câu 6:} Hệ bất phương trình bậc nhất hai ẩn là hệ gồm các bất phương trình bậc nhất hai ẩn $x$ và $y$ (bậc của mỗi biến đều là bậc nhất, không có số hạng bậc hai, tích $xy$ hay biến thứ ba). Hệ $\begin{cases} 3x - y + 2 \ge 0 \\ 2x + 5y - 7 < 0 \end{cases}$ thỏa mãn định nghĩa. Chọn \textbf{B}.
    \item \textbf{Câu 7:} Thay tọa độ điểm $N(4; 1)$ vào bất phương trình $3x - 2y \ge 6$:
    Ta có $3(4) - 2(1) = 12 - 2 = 10 \ge 6$ (luôn đúng). Do đó $N(4; 1)$ thuộc miền nghiệm. Chọn \textbf{C}.
    \item \textbf{Câu 8:} Bất phương trình $4x - 5y + 1 \ge 0$ có dạng $ax + by + c \ge 0$ với $a = 4 \ne 0, b = -5 \ne 0$, là bất phương trình bậc nhất hai ẩn. Chọn \textbf{D}.
    \item \textbf{Câu 9:} Theo công thức tính độ dài đường trung tuyến trong tam giác:
    $m_a^2 = \dfrac{2(b^2 + c^2) - a^2}{4}$. Chọn \textbf{A}.
    \item \textbf{Câu 10:} Ta có $\sin 150^\circ = \sin(180^\circ - 30^\circ) = \sin 30^\circ = \dfrac{1}{2}$. Chọn \textbf{B}.
    \item \textbf{Câu 11:} Với hai góc bù nhau $\alpha$ và $180^\circ - \alpha$, ta có mối quan hệ:
    $\sin(180^\circ - \alpha) = \sin \alpha$, $\cos(180^\circ - \alpha) = -\cos \alpha$, $\tan(180^\circ - \alpha) = -\tan \alpha$. Khẳng định đúng là $\cos(180^\circ - \alpha) = -\cos \alpha$. Chọn \textbf{C}.
    \item \textbf{Câu 12:} Với góc tù $90^\circ < \alpha < 180^\circ$, điểm ngọn của cung lượng giác nằm ở góc phần tư thứ II trên nửa đường tròn đơn vị, khi đó hoành độ âm ($\cos \alpha < 0$) và tung độ dương ($\sin \alpha > 0$). Suy ra $\tan \alpha = \dfrac{\sin \alpha}{\cos \alpha} < 0$ và $\cot \alpha = \dfrac{\cos \alpha}{\sin \alpha} < 0$. Do đó khẳng định $\cot \alpha < 0$ là đúng. Chọn \textbf{D}.
\end{itemize}

\vspace{0.2cm}
\textbf{PHẦN II. Trắc nghiệm đúng/sai}
\begin{itemize}[leftmargin=0.6cm, itemsep=4pt]
    \item \textbf{Câu 13:} Cho $A = [-4; 3)$ và $B = (0; 5]$.
    \begin{enumerate}
        \item \textbf{Đúng.} Vì $-4 \le 0 < 3$ nên $0 \in A$; mặt khác điều kiện để thuộc $B$ là $0 < x \le 5$, số $0$ không thỏa mãn nên $0 \notin B$.
        \item \textbf{Đúng.} Theo định nghĩa nửa khoảng, $B = (0; 5] = \{x \in \mathbb{R} \mid 0 < x \le 5\}$.
        \item \textbf{Đúng.} $A \cap B = [-4; 3) \cap (0; 5] = (0; 3)$.
        \item \textbf{Sai.} Vì $0 \in A$ và $0 \notin B$ nên phần tử $0 \in (A \setminus B)$. Do đó mút $0$ phải lấy dấu ngoặc vuông: $A \setminus B = [-4; 0]$. Khẳng định nửa khoảng $[-4; 0)$ là sai.
    \end{enumerate}
    \item \textbf{Câu 14:} Cho góc nhọn $\alpha$ ($0^\circ < \alpha < 90^\circ$) và $\beta = 180^\circ - \alpha$.
    \begin{enumerate}
        \item \textbf{Đúng.} Vì $\alpha + \beta = \alpha + (180^\circ - \alpha) = 180^\circ$, theo định nghĩa $\alpha$ và $\beta$ là hai góc bù nhau.
        \item \textbf{Đúng.} Theo tính chất lượng giác của hai góc bù nhau, $\sin(180^\circ - \alpha) = \sin \alpha \implies \sin \beta = \sin \alpha$.
        \item \textbf{Đúng.} Vì hai góc $\alpha$ và $90^\circ - \alpha$ phụ nhau nên $\sin(90^\circ - \alpha) = \cos \alpha$. Do đó:
        $M = \sin^2 \alpha + \cos^2 \alpha = 1$.
        \item \textbf{Sai.} Khi $\alpha = 45^\circ$, ta có $\beta = 180^\circ - 45^\circ = 135^\circ$.
        Suy ra $\cos \alpha = \cos 45^\circ = \dfrac{\sqrt{2}}{2}$, $\cos \beta = \cos 135^\circ = -\dfrac{\sqrt{2}}{2} \implies \cos \alpha + \cos \beta = 0$.
        Đồng thời $\tan \alpha = \tan 45^\circ = 1$, $\cot \beta = \cot 135^\circ = -1 \implies \tan \alpha + \cot \beta = 0$.
        Vậy $N = 0 + 0 = 0 \ne 2$. Khẳng định $N = 2$ là sai.
    \end{enumerate}
\end{itemize}

\vspace{0.2cm}
\textbf{PHẦN III. Trắc nghiệm trả lời ngắn}
\begin{itemize}[leftmargin=0.6cm, itemsep=4pt]
    \item \textbf{Câu 15:} Giao của hai tập hợp $A = [-3; 7)$ và $B = (1; 10]$ là:
    $C = A \cap B = [-3; 7) \cap (1; 10] = (1; 7)$.
    Các số nguyên thuộc khoảng $(1; 7)$ thỏa mãn $1 < x < 7$ là: $\{2; 3; 4; 5; 6\}$.
    Vậy tập hợp $C$ có tất cả $5$ phần tử nguyên. \textbf{Đáp số: 5}.
    \item \textbf{Câu 16:} Xét biểu thức $f(x, y) = 2x + 3y$.
    Vì $f(x, y)$ là biểu thức bậc nhất hai ẩn và miền tam giác $ABC$ là một miền đa giác lồi đóng có các đỉnh là $A(1; 2)$, $B(4; 1)$, $C(2; 5)$, nên giá trị lớn nhất của $f(x, y)$ trên miền tam giác đạt tại một trong ba đỉnh:
    \begin{itemize}
        \item Tại $A(1; 2)$: $f(1, 2) = 2(1) + 3(2) = 8$.
        \item Tại $B(4; 1)$: $f(4, 1) = 2(4) + 3(1) = 11$.
        \item Tại $C(2; 5)$: $f(2, 5) = 2(2) + 3(5) = 19$.
    \end{itemize}
    Suy ra $\max_{(x, y) \in \Delta ABC} (2x + 3y) = 19$.
    Để bất phương trình $2x + 3y \le m$ nghiệm đúng với mọi điểm $(x; y)$ thuộc miền tam giác $ABC$, điều kiện cần và đủ là:
    $m \ge \max_{(x, y) \in \Delta ABC} (2x + 3y) = 19$.
    Vậy giá trị nguyên nhỏ nhất của $m$ là $19$. \textbf{Đáp số: 19}.
    \item \textbf{Câu 17:} Theo định lý côsin trong tam giác $ABC$, ta có:
    $a^2 = b^2 + c^2 - 2bc \cos A \iff b^2 + c^2 - a^2 = 2bc \cos A$.
    Thay vào hệ thức đã cho:
    $2bc \cos A = \sqrt{3}bc \iff \cos A = \dfrac{\sqrt{3}}{2}$ (do $b, c > 0$).
    Vì trong tam giác $0^\circ < A < 180^\circ$ nên $A = 30^\circ$. \textbf{Đáp số: 30}.
    \item \textbf{Câu 18:} Xét phương trình $(x - 2)(x^2 - 5x + 6) = 0 \iff (x - 2)^2(x - 3) = 0 \iff x = 2$ hoặc $x = 3$.
    Do $2, 3 \in \mathbb{N}$ nên tập hợp $A = \{2; 3\}$.
    Tập hợp $B = \{x \in \mathbb{Z} \mid -1 \le x \le 4\} = \{-1; 0; 1; 2; 3; 4\}$.
    Khi đó hiệu của hai tập hợp là:
    $B \setminus A = \{-1; 0; 1; 4\}$.
    Số phần tử của tập hợp $B \setminus A$ là $4$. \textbf{Đáp số: 4}.
\end{itemize}

\vspace{0.2cm}
\textbf{PHẦN IV. Tự luận}
\begin{itemize}[leftmargin=0.6cm, itemsep=6pt]
    \item \textbf{Câu 19:}
    \begin{enumerate}
        \item Áp dụng định lý côsin trong tam giác $ABC$:
        \[
        BC^2 = AB^2 + AC^2 - 2AB \cdot AC \cdot \cos \widehat{BAC}
        \]
        Thay các giá trị $AB = 50$, $AC = 80$, $\widehat{BAC} = 60^\circ$ vào công thức:
        \[
        BC^2 = 50^2 + 80^2 - 2 \cdot 50 \cdot 80 \cdot \cos 60^\circ = 2500 + 6400 - 8000 \cdot \dfrac{1}{2} = 8900 - 4000 = 4900.
        \]
        Suy ra $BC = \sqrt{4900} = 70\text{ m}$.
        \item Diện tích khu đất tam giác $ABC$ là:
        \[
        S = \dfrac{1}{2} AB \cdot AC \cdot \sin \widehat{BAC} = \dfrac{1}{2} \cdot 50 \cdot 80 \cdot \sin 60^\circ = 2000 \cdot \dfrac{\sqrt{3}}{2} = 1000\sqrt{3}\text{ m}^2.
        \]
        \item Nửa chu vi của tam giác $ABC$ là:
        \[
        p = \dfrac{AB + BC + CA}{2} = \dfrac{50 + 70 + 80}{2} = \dfrac{200}{2} = 100\text{ m}.
        \]
        Bán kính $r$ của hồ nước tròn nội tiếp tam giác $ABC$ là:
        \[
        r = \dfrac{S}{p} = \dfrac{1000\sqrt{3}}{100} = 10\sqrt{3}\text{ m} \approx 17{,}3\text{ m}.
        \]
    \end{enumerate}

    \item \textbf{Câu 20:}
    Gọi $A, B, C$ lần lượt là tập hợp các học sinh đăng ký thi đấu môn Bóng đá, Cầu lông và Bóng rổ.
    Theo giả thiết, ta có:
    $|A| = 25$, $|B| = 20$, $|C| = 15$;
    $|A \cap B| = 10$, $|B \cap C| = 7$, $|C \cap A| = 8$;
    $|A \cap B \cap C| = 4$.
    \begin{enumerate}
        \item Số học sinh đăng ký tham gia ít nhất một môn thể thao là số phần tử của tập hợp $A \cup B \cup C$. Áp dụng công thức bao hàm -- loại trừ:
        \[
        |A \cup B \cup C| = |A| + |B| + |C| - (|A \cap B| + |B \cap C| + |C \cap A|) + |A \cap B \cap C|
        \]
        \[
        = 25 + 20 + 15 - (10 + 7 + 8) + 4 = 60 - 25 + 4 = 39\text{ học sinh}.
        \]
        \item Số học sinh trong lớp 10A không tham gia môn thể thao nào là:
        \[
        45 - |A \cup B \cup C| = 45 - 39 = 6\text{ học sinh}.
        \]
        \item Số học sinh chỉ đăng ký tham gia đúng một môn thể thao được tính bằng:
        \begin{itemize}
            \item Chỉ tham gia Bóng đá: $|A| - (|A \cap B| + |A \cap C| - |A \cap B \cap C|) = 25 - (10 + 8 - 4) = 25 - 14 = 11$.
            \item Chỉ tham gia Cầu lông: $|B| - (|A \cap B| + |B \cap C| - |A \cap B \cap C|) = 20 - (10 + 7 - 4) = 20 - 13 = 7$.
            \item Chỉ tham gia Bóng rổ: $|C| - (|A \cap C| + |B \cap C| - |A \cap B \cap C|) = 15 - (8 + 7 - 4) = 15 - 11 = 4$.
        \end{itemize}
        Tổng số học sinh chỉ tham gia đúng một môn là: $11 + 7 + 4 = 22$ học sinh.
    \end{enumerate}

    \item \textbf{Câu 21:}
    Gọi $x, y$ lần lượt là số chiếc đèn lồng loại A và số chiếc hộp gỗ loại B mà cơ sở sản xuất trong một tuần ($x, y \in \mathbb{N}$).
    \begin{itemize}
        \item Thời gian gia công trong một tuần là $2x + y$ giờ, không vượt quá $100$ giờ: $2x + y \le 100$.
        \item Thời gian hoàn thiện trong một tuần là $x + 2y$ giờ, không vượt quá $80$ giờ: $x + 2y \le 80$.
        \item Điều kiện không âm: $x \ge 0, y \ge 0$.
    \end{itemize}
    Ta có hệ bất phương trình ràng buộc:
    \[
    \begin{cases}
    2x + y \le 100 \\
    x + 2y \le 80 \\
    x \ge 0 \\
    y \ge 0
    \end{cases}
    \]
    Miền nghiệm của hệ là miền tứ giác $OABC$ (kể cả biên) trên mặt phẳng tọa độ $Oxy$ với:
    \begin{itemize}
        \item $O(0; 0)$ là gốc tọa độ;
        \item $A(50; 0)$ là giao điểm của đường thẳng $2x + y = 100$ với trục hoành;
        \item $C(0; 40)$ là giao điểm của đường thẳng $x + 2y = 80$ với trục tung;
        \item $B(40; 20)$ là tọa độ giao điểm của hai đường thẳng $2x + y = 100$ và $x + 2y = 80$.
    \end{itemize}
    Tổng lợi nhuận thu được trong tuần (đơn vị: nghìn đồng) là:
    \[
    F(x, y) = 120x + 100y
    \]
    Vì $F(x, y)$ là biểu thức bậc nhất hai ẩn, giá trị lớn nhất của $F(x, y)$ trên miền đa giác đạt tại một trong các đỉnh:
    \begin{itemize}
        \item Tại $O(0; 0)$: $F(0, 0) = 0$.
        \item Tại $A(50; 0)$: $F(50, 0) = 120(50) + 100(0) = 6000$ (nghìn đồng).
        \item Tại $C(0; 40)$: $F(0, 40) = 120(0) + 100(40) = 4000$ (nghìn đồng).
        \item Tại $B(40; 20)$: $F(40, 20) = 120(40) + 100(20) = 4800 + 2000 = 6800$ (nghìn đồng).
    \end{itemize}
    So sánh giá trị tại các đỉnh, $F(x, y)$ đạt giá trị lớn nhất bằng $6800$ nghìn đồng khi $x = 40$ và $y = 20$.
    Vậy mỗi tuần cơ sở cần sản xuất \textbf{40 chiếc đèn lồng loại A} và \textbf{20 chiếc hộp gỗ loại B} để thu được lợi nhuận lớn nhất là \textbf{6 800 000 đồng}.

    \item \textbf{Câu 22:}
    Gọi $A$ là tập hợp học sinh tham gia CLB Tin học, $B$ là tập hợp học sinh tham gia CLB Tiếng Anh.
    Theo bài ra, ta có $|A| = 32$, $|B| = 28$, và số bạn tham gia cả hai CLB là $|A \cap B| = 18$.
    \begin{enumerate}
        \item Số học sinh tham gia ít nhất một trong hai câu lạc bộ là:
        \[
        |A \cup B| = |A| + |B| - |A \cap B| = 32 + 28 - 18 = 42\text{ bạn}.
        \]
        \item Số học sinh không tham gia cả hai câu lạc bộ trên là:
        \[
        50 - |A \cup B| = 50 - 42 = 8\text{ bạn}.
        \]
        \item Số học sinh chỉ tham gia Câu lạc bộ Tin học mà không tham gia Câu lạc bộ Tiếng Anh là:
        \[
        |A \setminus B| = |A| - |A \cap B| = 32 - 18 = 14\text{ bạn}.
        \]
    \end{enumerate}
\end{itemize}

\end{document}
